\section{Tabu Search}
Nelle slide venivano proposte tre metaeuristiche: Simulated Annealing (SA, ricottura simulata), Tabu Search (TS), Genetic Algorithms (ispirati alla natura).
Riporto una descrizione generale di SA data da MATLAB:
\begin{quote}
	Simulated annealing is a method for solving unconstrained and bound-constrained optimization problems. The method models the physical process of heating a material and then slowly lowering the temperature to decrease defects, thus minimizing the system energy.
	
	At each iteration of the simulated annealing algorithm, a new point is randomly generated. The distance of the new point from the current point, or the extent of the search, is based on a probability distribution with a scale proportional to the temperature. The algorithm accepts all new points that lower the objective, \textbf{but also, with a certain probability, points that raise the objective}. By accepting points that raise the objective, \textbf{the algorithm avoids being trapped in local minima, and is able to explore globally for more possible solutions}. An annealing schedule is selected to systematically decrease the temperature as the algorithm proceeds. As the temperature decreases, the algorithm reduces the extent of its search to converge to a minimum
\end{quote}
Questa definizione è utile soltanto per introdurre qualche concetto utile alla Tabu Search. 

Diamo due definizioni:
\begin{itemize}
	\item \textbf{Mossa}: operazione eseguita per definire il vicinato di una data soluzione;
	\item \textbf{Mossa tabù}: mossa proibita.
\end{itemize}
Evito di rieseguire/cancellare gli effetti di una mossa appena fatta rendendola tabù, questo per evitare di rimanere bloccati in ottimi locali. Le caratteristiche principali della TS sono:
\begin{itemize}
	\item Accettazione di mosse che deteriorano il valore della f.o;
	\item Costruzione di una \textbf{lista tabù} che contiene tutte le mosse proibite (per evitare di ritornare in un minimo locale);
	\item Utilizzo di due tipi di memorie: \textbf{Short Term Memory, Long Term Memory};
	\item Due fasi: \textbf{Intensificazione} e \textbf{Diversificazione};
	\item \textbf{Aspiration Criteria}: eventuali condizioni che permettono di eseguire una mossa tabù, questo per evitare restrizioni troppo forti.
\end{itemize}
La differenza tra i due tipi di memoria è fondamentale, in quanto entrambe hanno la loro specifica strategia: questo permette la modifica del vicinato $\mathcal{N}(x)$ della soluzione corrente $x$ e il vicinato modificato è il risultato di una storia selettiva degli stati incontrati durante il processo di ricerca (quella a breve termine ha una tabu list per definire elementi del vicinato esclusi da quello modificato, mentre in quella a lungo termine posso includere soluzioni che non appartengono al vicinato originale). TS è un metodo con vicinato dinamico.

L'algoritmo base con sola Short Term Memory (STM, e l'altra LTM);
\begin{enumerate}
	\item Identifica una soluzione ammissibile $\tilde{x}$: imposta $x:=\tilde{x}$ (con $x$ soluzione attuale) e inizializza la Tabu List $TL:=\emptyset$;
	\item Genera vicinato $\mathcal{N}(x)$;
	\item Trova soluzione $y\in \mathcal{N}(x)$ tale che:
	\begin{itemize}
		\item $f(y)$ è minimo;
		\item $y\neq x$;
		\item la soluzione non mosse tabù.
	\end{itemize}
	\item Imposto $x:=y$;
	\item Aggiorno TL mettendo la mossa che cambia $x$ in $y$;
	\item Se non è soddisfatto il criterio di stop (n iterazioni massimo o da ultimo miglioramento, limite tempo), torna allo step 2 altrimenti stop.
\end{enumerate}
Le memorie tengono tracce delle soluzioni secondo:
\begin{itemize}
	\item Recency: STM più usata, tiene traccia delle soluzioni con attributi recentemente cambiati, usata per la Tabu List (\textbf{Tabu Tenure} per indicare per quante iterazioni una mossa rimane tabù);
	\item Frequency;
	\item Quality;
	\item Influence: quanto una caratteristica specifica cambia non solo il valore della soluzione, ma anche la struttura stessa.
\end{itemize}
Si aggiungono anche \textbf{Aspiration Plus}, soglia per qualità affinché una soluzione venga tenuta, e \textbf{Elite list}, lista di soluzioni migliori incontrate durante il processo di ricerca.

Da ricordarsi: la TS NECESSITA una soluzione iniziale dalla quale partire (solitamente ottenuta con un algoritmo greedy).

La CEM (\textbf{Critical Event Memory}) è una strategia (non una memoria) che può essere usata per fermare il processo di ricerca e riniziare da una soluzione differente e anche proveddere diversificazione. Registra eventi critici, ad esempio trovare nuove soluzioni iniziali o una mossa che migliora la soluzione precedente.

Abbiamo definito la tabu tenure come durata $j$ dello stato tabu-attivo per una mossa, un valore basso permette un'intensificazione, un valore alto permette la diversificazione (obbliga a spostarsi). Si usa un range di valori solitamente, non uno fisso, e può variare o in modo random (dopo $t_{max}$ iterazioni o ogni volta che un attributo viene aggiunto alla TL) o in modo sistematico (quando succede qualcosa).

Una mossa ad alta influenza non per forza migliora la f.o, ma cambia la struttura della soluzione, permettendo di visitarne di nuove.

Vediamo ora invece degli aspiration criteria:
\begin{itemize}
	\item Default: tutte mosse tabù, devo sceglierne una, quella meno tabù;
	\item Obiettivo: seleziono mossa che trova soluzione migliore di quella visitata fino a quel momento;
	\item Direzione di ricerca: un attributo può essere aggiunto/tolto da una soluzione se la direzione della ricerca non cambia;
	\item Influenza: faccio mossa tabù poco influente se ne ho appena fatta una altamente influente.
\end{itemize}

Infine, la LTM tiene misura di \textbf{Transition}, quanto una caratteristica viene modificata nel tempo dall'algoritmo (più alto è più è probabile che la caratteristica non sia di importanza per la struttura della soluzione), e di \textbf{Residence}, quanto una caratteristica appare nelle soluzioni visitate dall'algoritmo (se compare tanto, o è perché la caratteristica è un buon candidato o perché serve diversificare).
In questo caso, l'intensificazione usa lista di soluzioni d'elite, mentre la diversificazione sfrutta i dati salvati per fare dei cambiamenti alla struttura della soluzione.
